\documentclass[11pt]{beamer}
\usetheme{Madrid}
\usepackage[utf8]{inputenc}
\usepackage{amsmath}
\usepackage{amsfonts}
\usepackage{amssymb}
\usepackage{graphicx}
\usepackage{xcolor}
\usepackage{tikz}
\usepackage{geometry}

\title{EB1A \& EB2-NIW Citation Strategy}
\subtitle{Govt Citations • EDU Contributions • Major Significance • OCMS}
\author{Based on Strategic Research Methodology}
\date{}

\setbeamercovered{transparent}
\setbeamertemplate{navigation symbols}{}
\setbeamerfont{itemize/enumerate body}{size=\small}
\setbeamerfont{itemize/enumerate subbody}{size=\scriptsize}
\setbeamerfont{block body}{size=\small}
\setbeamerfont{block title}{size=\normalsize}

\begin{document}
	
	% TITLE SLIDE
	\begin{frame}
		\titlepage
		\centering
		\small \textit{How to Get .gov, .edu, and Journal Citations for Immigration Success}
	\end{frame}
	
	% ============================================================
	% SECTION 1: INTRODUCTION TO EB1A/EB2-NIW CITATION STRATEGY
	% ============================================================
	
	\begin{frame}{EB1A Prong 1: What USCIS Looks For}
		\small
		\begin{block}{Original Contributions of Major Significance}
			\textbf{Answer:} USCIS requires evidence that your work has had major significance and sustained acclaim in your field.
			\vspace{0.1cm}
			\textbf{Key evidence types:}
			\begin{itemize}
				\item Citations from .gov domains (federal agencies)
				\item Citations from .edu domains (universities)
				\item High-impact journal citations (Scopus, Web of Science)
				\item Evidence your work influenced policy, regulation, or practice
				\item Expert opinions confirming your contributions
			\end{itemize}
			\textcolor{blue}{Volume alone is not enough - quality and significance matter most.}
		\end{block}
	\end{frame}
	
	\begin{frame}{EB2-NIW: What USCIS Requires}
		\small
		\begin{block}{National Interest Waiver Criteria}
			\textbf{Answer:} EB2-NIW requires demonstrating three things:
			\vspace{0.1cm}
			\begin{enumerate}
				\item Your proposed endeavor has substantial merit and national importance
				\item You are well-positioned to advance the proposed endeavor
				\item On balance, it would be beneficial to the U.S. to waive the job offer and labor certification
			\end{enumerate}
			\textcolor{blue}{Government citations and policy influence directly prove "national importance" and "substantial merit."}
		\end{block}
	\end{frame}
	
	\begin{frame}{Why .gov Citations Are So Powerful for Immigration}
		\small
		\begin{block}{Government Domain Citations - Highest Credibility}
			\textbf{Answer:} .gov citations are considered high-value evidence because:
			\vspace{0.1cm}
			\begin{itemize}
				\item Federal agencies do not cite irrelevant or low-quality work
				\item A .gov citation proves your work influenced U.S. government policy
				\item USCIS officers recognize .gov domains as authoritative sources
				\item Verifiable - the citation can be checked online instantly
				\item Demonstrates "national importance" for NIW cases
				\item Proves "major significance" for EB1A petitions
			\end{itemize}
			\textcolor{red}{A single .gov citation can be worth dozens of regular journal citations.}
		\end{block}
	\end{frame}
	
	\begin{frame}{Understanding Citations vs. Indexing vs. Commenting}
		\small
		\begin{block}{Three Different Types of Evidence}
			\textbf{Answer:} Each type serves a different purpose:
			\vspace{0.1cm}
			\textbf{Citations:} Your work directly referenced by another author; found in Google Scholar, Scopus; most valuable for EB1A.
			\vspace{0.1cm}
			\textbf{Indexing:} Your work stored in a database (ERIC, PubMed, Science.gov); shows discoverability; good for NIW.
			\vspace{0.1cm}
			\textbf{Commenting:} Submit feedback to Regulations.gov; creates official public record; excellent for policy influence.
		\end{block}
	\end{frame}
	
	\begin{frame}{How Many Citations Are Enough for EB1A?}
		\small
		\begin{block}{Citation Targets - What Actually Works}
			\textbf{Answer:} No fixed number, but realistic targets:
			\vspace{0.1cm}
			\textbf{Minimum (with strong other evidence):}
			\begin{itemize}
				\item 50-100 total citations
				\item At least 5-10 .gov or policy citations
			\end{itemize}
			\textbf{Competitive (strong EB1A case):}
			\begin{itemize}
				\item 200-500+ total citations
				\item Multiple .gov and .edu citations
				\item Peer review for prestigious journals
			\end{itemize}
			\textcolor{blue}{Quality matters more than quantity - one policy citation can be worth 50 journal citations.}
		\end{block}
	\end{frame}
	
	\begin{frame}{Setting Realistic Citation Targets by Field}
		\small
		\begin{block}{Field-Specific Expectations}
			\textbf{Answer:} Citation norms vary by discipline:
			\vspace{0.1cm}
			\textbf{Social Sciences \& Humanities:} 50-150 citations = strong; focus on .gov citations.
			\vspace{0.1cm}
			\textbf{STEM Fields:} 200-500 citations = competitive; 500+ = extraordinary.
			\vspace{0.1cm}
			\textbf{Computer Science \& Engineering:} Conference papers cited less; GitHub stars matter too.
		\end{block}
	\end{frame}
	
	% ============================================================
	% SECTION 2: GOVERNMENT CITATION STRATEGY (.GOV DOMAINS)
	% ============================================================
	
	\begin{frame}{How to Get .gov Citations: The Complete Strategy}
		\small
		\begin{block}{Six Proven Methods for Government Citations}
			\textbf{Answer:} Six reliable ways to get cited by government agencies:
			\vspace{0.1cm}
			\begin{enumerate}
				\item \textbf{Regulations.gov Commenting}: Submit research-backed comments
				\item \textbf{Federal Agency Reports}: Get cited in BLS, Fed, DOE reports
				\item \textbf{CLEAR Database}: Create datasets for DOL's CLEAR system
				\item \textbf{Government RFI Responses}: Respond to agency RFIs
				\item \textbf{Federal Contract Work}: Work on government-funded research
				\item \textbf{Collaboration with Federal Researchers}: Co-author with government scientists
			\end{enumerate}
			\textcolor{red}{Regulations.gov commenting is the most accessible method.}
		\end{block}
	\end{frame}
	
	\begin{frame}{Regulations.gov: How to Submit Effective Comments}
		\small
		\begin{block}{Step-by-Step Commenting Strategy}
			\textbf{Step 1 - Find Relevant RFIs:}
			\begin{itemize}
				\item Visit Regulations.gov and search your topic
				\item Look for "Request for Information" (RFI) or "Notice of Proposed Rulemaking" (NPRM)
			\end{itemize}
			\textbf{Step 2 - Write Research-Backed Comments:}
			\begin{itemize}
				\item Cite your own published research
				\item Reference peer-reviewed literature
				\item Provide data-driven recommendations
				\item Keep professional tone (500-2000 words)
			\end{itemize}
		\end{block}
	\end{frame}
	
	\begin{frame}{Regulations.gov Commenting: Formatting for Success}
		\small
		\begin{block}{How to Format Your Comments}
			\textbf{Required elements:}
			\begin{itemize}
				\item Docket number (e.g., "Docket No. EPA-HQ-OAR-2022-1234")
				\item Your name and affiliation
				\item Clear subject line
				\item Executive summary (1-2 paragraphs)
				\item Detailed analysis with citations
				\item Specific recommendations
			\end{itemize}
			\textbf{Pro tip:} Use LaTeX to format professionally before uploading as PDF.
			\textcolor{blue}{Comments become official public record and can be cited by others.}
		\end{block}
	\end{frame}
	
	\begin{frame}{How to Track When Your Comment Gets Cited}
		\small
		\begin{block}{Monitoring Your Policy Impact}
			\textbf{What to watch for:}
			\begin{itemize}
				\item Agency final rules may reference comments received
				\item Federal agency reports may cite your research
				\item Other commenters may reference your submission
				\item Google Scholar may index the final rule (with your citation!)
			\end{itemize}
			\textbf{Documentation strategy:}
			\begin{itemize}
				\item Save docket number and confirmation email
				\item Screenshot your comment on the public portal
				\item Download PDFs for your portfolio
			\end{itemize}
		\end{block}
	\end{frame}
	
	\begin{frame}{CLEAR DOL: Creating Datasets for Department of Labor}
		\small
		\begin{block}{CLEAR Database - What It Is}
			\textbf{Answer:} CLEAR (Clearinghouse for Labor Evaluation and Research) is DOL's evidence database.
			\vspace{0.1cm}
			\textbf{How to get indexed:}
			\begin{itemize}
				\item Publish rigorous research on labor topics (workforce, jobs, training)
				\item Ensure strong methodology (RCT, quasi-experimental)
				\item Submit through CLEAR online portal
				\item DOL reviews and indexes quality studies
			\end{itemize}
			\textbf{CLEAR indexing proves:} Research meets federal quality standards and influences labor policy.
		\end{block}
	\end{frame}
	
	\begin{frame}{Bureau of Labor Statistics (BLS) Citations}
		\small
		\begin{block}{How to Get Cited by BLS}
			\textbf{Methods:}
			\begin{itemize}
				\item Publish on labor economics, workforce trends, employment statistics
				\item Submit comments to BLS RFIs
				\item Collaborate with BLS researchers
				\item Get referenced in BLS Monthly Labor Review
			\end{itemize}
			\textbf{How to check if BLS cited you:}
			\begin{itemize}
				\item Search BLS publications for your name
				\item Use Google Scholar with "site:bls.gov" + your name
				\item Check Crossref XML for DOI resolution
			\end{itemize}
			\textcolor{blue}{BLS citations prove your work influences national workforce policy.}
		\end{block}
	\end{frame}
	
	\begin{frame}{Federal Reserve Citations: Getting Cited by the Fed}
		\small
		\begin{block}{How to Get Cited by Federal Reserve}
			\textbf{Strategies:}
			\begin{itemize}
				\item Publish on monetary policy, banking, financial stability
				\item Submit comments to Fed RFIs (payment systems, crypto regulation)
				\item Collaborate with Fed economists
				\item Get cited in FEDS Working Papers
			\end{itemize}
			\textbf{Key Fed publications to target:}
			\begin{itemize}
				\item FEDS Working Papers
				\item FedNotes
				\item Beige Books
				\item Consumer Compliance Outlook
			\end{itemize}
		\end{block}
	\end{frame}
	
	\begin{frame}{Science.gov Indexing: Federal Repositories}
		\small
		\begin{block}{Science.gov - Gateway to Federal Science}
			\textbf{How to get indexed:}
			\begin{itemize}
				\item Publish in journals indexed by PubMed, ERIC, or DOAJ
				\item Submit to federal agency repositories (NASA ADS, DOE OSTI)
				\item Ensure your work has a DOI
				\item Use clear, searchable titles and abstracts
			\end{itemize}
			\textbf{Pathways to Science.gov:}
			\begin{itemize}
				\item NASA ADS $\rightarrow$ Science.gov
				\item Harvard ADS $\rightarrow$ Science.gov
				\item ERIC $\rightarrow$ Science.gov
				\item PubMed $\rightarrow$ Science.gov
			\end{itemize}
			\textcolor{blue}{Indexing proves federal recognition of your research.}
		\end{block}
	\end{frame}
	
	% ============================================================
	% SECTION 3: ERIC.GOV AND EDUCATIONAL IMPACT
	% ============================================================
	
	\begin{frame}{ERIC.gov: What Is It and Why It Matters}
		\small
		\begin{block}{Education Resources Information Center (ERIC)}
			\textbf{Answer:} ERIC is the U.S. Department of Education's digital library.
			\vspace{0.1cm}
			\textbf{Why ERIC indexing is valuable:}
			\begin{itemize}
				\item .gov domain authority
				\item Work becomes searchable by educators and policymakers
				\item Demonstrates educational impact for NIW
				\item Can be cited by .edu domains
				\item Free access increases citation potential
			\end{itemize}
		\end{block}
	\end{frame}
	
	\begin{frame}{How to Get Your Work Indexed in ERIC}
		\small
		\begin{block}{ERIC Indexing Requirements}
			\textbf{Pathways to ERIC indexing:}
			\begin{enumerate}
				\item Publish in ERIC-approved journals (many education journals)
				\item Submit directly to ERIC (unpublished research, reports, dissertations)
				\item Conference proceedings (some education conferences)
				\item Open peer review journals (if ERIC-approved)
			\end{enumerate}
			\textbf{Journal characteristics:} Peer-reviewed, English, education focus, has ISSN.
		\end{block}
	\end{frame}
	
	\begin{frame}{Open Peer Review Journals with ERIC Indexing}
		\small
		\begin{block}{Best Journals for ERIC Indexing}
			\textbf{Mid-tier (easier acceptance):}
			\begin{itemize}
				\item Journal of Education and Learning (EduLearn)
				\item International Journal of Instruction
				\item Journal of Technology and Science Education
			\end{itemize}
			\textbf{Higher-tier (more citations):}
			\begin{itemize}
				\item Computers \& Education (Elsevier) - Scopus Q1
				\item Educational Technology Research and Development (Springer)
				\item British Journal of Educational Technology (Wiley)
			\end{itemize}
			\textcolor{blue}{Aim for mid-tier first, then higher-tier after establishing record.}
		\end{block}
	\end{frame}
	
	\begin{frame}{How ERIC Citations Help Your Immigration Case}
		\small
		\begin{block}{Demonstrating Educational Impact}
			\textbf{Evidence for EB1A/EB2-NIW:}
			\begin{itemize}
				\item ERIC is a .gov domain - automatic credibility
				\item Proves research reaches educators nationwide
				\item Demonstrates "substantial merit" for NIW
				\item Counts as "professional recognition"
			\end{itemize}
			\textbf{How to document:}
			\begin{itemize}
				\item Save ERIC abstract page with .gov URL
				\item Track citations to your ERIC-indexed work
				\item Monitor download statistics
			\end{itemize}
		\end{block}
	\end{frame}
	
	% ============================================================
	% SECTION 4: HIGH-IMPACT PUBLICATION STRATEGY
	% ============================================================
	
	\begin{frame}{Journal Selection Checklist for Maximum Citations}
		\small
		\begin{block}{How to Identify High-Citation-Potential Journals}
			\textbf{Checklist:}
			\begin{itemize}
				\item [$\checkmark$] Indexed in Scopus, Web of Science, or DOAJ
				\item [$\checkmark$] Has Journal Impact Factor
				\item [$\checkmark$] Published by recognized publisher (Elsevier, Springer, Wiley, Taylor \& Francis)
				\item [$\checkmark$] No predatory red flags
				\item [$\checkmark$] Quick publication (2-6 months)
				\item [$\checkmark$] Open access option available
				\item [$\checkmark$] Has DOIs for all articles
				\item [$\checkmark$] Indexed in Google Scholar
			\end{itemize}
			\textcolor{red}{Avoid predatory journals - they cause RFEs.}
		\end{block}
	\end{frame}
	
	\begin{frame}{Research Topics with High Citation Velocity}
		\small
		\begin{block}{Topics Getting Cited Rapidly (2024-2026)}
			\textbf{Artificial Intelligence:}
			\begin{itemize}
				\item LLMs in various applications
				\item AI ethics, governance, regulation
				\item Generative AI in education and workforce
			\end{itemize}
			\textbf{Cybersecurity:}
			\begin{itemize}
				\item Zero-trust architecture
				\item AI security and adversarial ML
				\item Critical infrastructure protection
			\end{itemize}
			\textbf{FinTech \& Financial Regulation:}
			\begin{itemize}
				\item Cryptocurrency regulation
				\item Digital assets and stablecoins
				\item AI in banking compliance
			\end{itemize}
		\end{block}
	\end{frame}
	
	\begin{frame}{Cross-Disciplinary Research for Broader Impact}
		\small
		\begin{block}{How to Increase Citations Through Interdisciplinary Work}
			\textbf{High-impact combinations:}
			\begin{itemize}
				\item AI + Education (AI literacy, personalized learning)
				\item AI + Finance (algorithmic trading risk, fraud detection)
				\item AI + Cybersecurity (adversarial ML, threat detection)
				\item FinTech + Regulation (compliance automation)
			\end{itemize}
			\textbf{Why it works:} Reaches multiple communities, cited across fields, attractive to RFIs.
		\end{block}
	\end{frame}
	
	\begin{frame}{LaTeX Templates for Professional Paper Submission}
		\small
		\begin{block}{Using LaTeX for Academic Publishing}
			\textbf{Why use LaTeX:}
			\begin{itemize}
				\item Preferred by most STEM journals
				\item Automatic formatting to journal specs
				\item Handles references automatically (BibTeX)
				\item Free (Overleaf, TeX Live)
			\end{itemize}
			\textbf{Where to get templates:}
			\begin{itemize}
				\item Overleaf (overleaf.com/latex/templates)
				\item Journal websites (download .cls files)
				\item GitHub
			\end{itemize}
			\textcolor{blue}{Learn basic LaTeX - worth the investment.}
		\end{block}
	\end{frame}
	
	\begin{frame}{Zotero: Citation Management for Immigration Evidence}
		\small
		\begin{block}{Using Zotero to Track Your Citations}
			\textbf{Why Zotero for immigration:}
			\begin{itemize}
				\item Automatically pulls citation data from Google Scholar
				\item Creates formatted bibliographies
				\item Generates citation reports
				\item Stores PDFs
				\item Free and open source
			\end{itemize}
			\textbf{How to use:}
			\begin{itemize}
				\item Create "My Publications" collection
				\item Add all papers, preprints, comments
				\item Use "Create Citation Report" to show times cited
				\item Track who cites you (.gov, .edu, journals)
			\end{itemize}
		\end{block}
	\end{frame}
	
	\begin{frame}{Using AI (DeepSeek, ChatGPT) for Research Writing}
		\small
		\begin{block}{AI Tools to Accelerate Publication}
			\textbf{How to use AI responsibly:}
			\begin{itemize}
				\item Generate literature review summaries
				\item Edit and polish existing writing
				\item Create LaTeX code for tables/figures
				\item Brainstorm research questions
			\end{itemize}
			\textbf{Be careful:}
			\begin{itemize}
				\item Always verify AI-generated information
				\item Do not submit AI text without editing
				\item Many journals prohibit AI as co-author
				\item \textbf{Do not use for peer review} - violates policies
			\end{itemize}
		\end{block}
	\end{frame}
	
	% ============================================================
	% SECTION 5: PEER REVIEW OPPORTUNITIES
	% ============================================================
	
	\begin{frame}{How to Get Peer Review Opportunities}
		\small
		\begin{block}{Building Your Peer Review Portfolio}
			\textbf{Strategies:}
			\begin{enumerate}
				\item Start small: Review for new or lower-tier journals
				\item Contact journal editors: Send CV and expertise
				\item Complete reviewer profiles: Register on ScholarOne, Editorial Manager
				\item Publish first: Editors invite established authors
				\item Ask mentors: Senior colleagues can recommend you
			\end{enumerate}
			\textbf{Evidence to keep:} Invitation emails, review confirmations, certificates, editor thank-you letters.
		\end{block}
	\end{frame}
	
	\begin{frame}{Peer Review Certifications from Publishers}
		\small
		\begin{block}{Which Publishers Provide Reviewer Certificates}
			\textbf{Publishers with certification programs:}
			\begin{itemize}
				\item \textbf{Elsevier}: Web of Science Reviewer Recognition
				\item \textbf{Wiley}: Reviewer Certificates available
				\item \textbf{Springer}: Publons integration
				\item \textbf{Taylor \& Francis}: Reviewer recognition
				\item \textbf{MDPI}: Immediate certificate download
				\item \textbf{Frontiers}: Certificates with feedback
			\end{itemize}
			\textbf{How to collect:} Create Publons profile, connect ORCID, upload confirmations, download certificates.
		\end{block}
	\end{frame}
	
	\begin{frame}{How to Document Peer Review for USCIS}
		\small
		\begin{block}{Evidence Package for "Judging the Work of Others"}
			\textbf{Required evidence:}
			\begin{enumerate}
				\item Invitation to review (journal email)
				\item Confirmation of completed review
				\item Reviewer certificate (if available)
				\item Screenshot of acknowledgment page
				\item Editor's thank-you letter
			\end{enumerate}
			\textbf{What USCIS looks for:}
			\begin{itemize}
				\item Prestigious journals (impact factor)
				\item Multiple reviews (5+ minimum)
				\item Span of time (1-3 years)
				\item Invitations based on expertise
			\end{itemize}
		\end{block}
	\end{frame}
	
	\begin{frame}{Mixture of Experts: Non-Peer-Reviewed Long Papers}
		\small
		\begin{block}{Using Long-Form Articles to Build Authority}
			\textbf{What "Mixture of Experts" means:}
			\begin{itemize}
				\item Invited book chapters
				\item Encyclopedia entries
				\item Professional magazine articles (IEEE Spectrum)
				\item Industry white papers
				\item Policy briefs
				\item Technical reports
			\end{itemize}
			\textbf{How to use:} Showcase breadth, reach different audiences, generate organic citations.
		\end{block}
	\end{frame}
	
	% ============================================================
	% SECTION 6: .EDU CITATION STRATEGY
	% ============================================================
	
	\begin{frame}{How to Get Cited by US Universities (.edu Domains)}
		\small
		\begin{block}{Five Methods for .edu Citations}
			\begin{enumerate}
				\item Collaborate with university researchers - co-author papers
				\item Get cited in course materials - reach out to professors
				\item University press publications - submit to university journals
				\item Institutional repositories - deposit in Harvard DASH, MIT DSpace
				\item Present at university seminars - get referenced in proceedings
			\end{enumerate}
			\textbf{Most accessible:} Identify professors, send publications, offer materials, ask to be considered for course readings.
		\end{block}
	\end{frame}
	
	\begin{frame}{Tracking .edu Citations Using Google Scholar}
		\small
		\begin{block}{How to Find University Citations}
			\textbf{Step-by-step:}
			\begin{enumerate}
				\item Go to Google Scholar My Profile
				\item Click "Cited by" number for your work
				\item Look at URLs of citing works
				\item Filter for URLs containing ".edu"
				\item Save these citations to your portfolio
			\end{enumerate}
			\textbf{Advanced search:} "your name" site:.edu
			\vspace{0.2cm}
			\textcolor{blue}{Citations from Harvard, MIT, Stanford are especially valuable.}
		\end{block}
	\end{frame}
	
	% ============================================================
	% SECTION 7: NIH AND PUBMED CITATIONS
	% ============================================================
	
	\begin{frame}{PubMed.gov: Getting Cited in Medical Literature}
		\small
		\begin{block}{PubMed Indexing and Citation Strategy}
			\textbf{How to get PubMed citations:}
			\begin{itemize}
				\item Publish in PubMed-indexed journals
				\item Focus on NIH mission topics (public health, medicine, health policy)
				\item Collaborate with NIH-funded researchers
				\item Submit to medRxiv preprint server
			\end{itemize}
			\textbf{Why PubMed matters:}
			\begin{itemize}
				\item .gov domain authority
				\item Accessible to U.S. health policymakers
				\item Demonstrates health-related national importance
			\end{itemize}
		\end{block}
	\end{frame}
	
	% ============================================================
	% SECTION 8: PROMOTING YOUR RESEARCH
	% ============================================================
	
	\begin{frame}{Upload Your Paper Everywhere: Multi-Platform Strategy}
		\small
		\begin{block}{Maximum Visibility Through Multiple Platforms}
			\textbf{Where to upload:}
			\begin{itemize}
				\item \textbf{Preprint servers}: arXiv, SSRN, medRxiv, bioRxiv
				\item \textbf{Institutional repositories}: University digital archive
				\item \textbf{ResearchGate}, \textbf{Academia.edu}
				\item \textbf{HAL}: French open archive
				\item \textbf{MPRA}, \textbf{EconPapers} (economics)
				\item \textbf{EBSCO}: Commercial database
			\end{itemize}
			\textcolor{blue}{Each platform increases discoverability and citation potential.}
		\end{block}
	\end{frame}
	
	\begin{frame}{EBSCO, ProQuest, and HAL: Getting Indexed}
		\small
		\begin{block}{Commercial and International Indexing}
			\textbf{EBSCOhost:} Used by 1000s of university libraries; publish in EBSCO-indexed journals.
			\vspace{0.1cm}
			\textbf{ProQuest:} Dissertations, theses, conference proceedings; submit through university library.
			\vspace{0.1cm}
			\textbf{HAL:} French government-supported; indexed by Google Scholar; accepts English submissions.
		\end{block}
	\end{frame}
	
	\begin{frame}{US Commerce Library Indexing Through EBSCO}
		\small
		\begin{block}{Getting Your Work in the U.S. Commerce Research Library}
			\textbf{How to get indexed:}
			\begin{itemize}
				\item Publish in EBSCO's Business Source Complete journals
				\item Focus on trade, economics, business, technology
				\item Ensure journal is in Commerce Library holdings
			\end{itemize}
			\textbf{Why this matters:}
			\begin{itemize}
				\item Direct .gov indexing
				\item Accessible to U.S. Department of Commerce policymakers
				\item Verifiable by USCIS through .gov website
			\end{itemize}
		\end{block}
	\end{frame}
	
	% ============================================================
	% SECTION 9: SPECIALIZED DOMAINS (CYBERSECURITY, AI)
	% ============================================================
	
	\begin{frame}{Cybersecurity as the Second Most Important Topic}
		\small
		\begin{block}{Why Cybersecurity Is Highly Valued}
			\textbf{Why cybersecurity citations are powerful:}
			\begin{itemize}
				\item Direct national importance for NIW
				\item Federal agencies cite security research (NIST, CISA, NSA)
				\item Executive Orders prioritize cybersecurity
				\item Shortage of U.S. cybersecurity experts
			\end{itemize}
			\textbf{Key government resources:}
			\begin{itemize}
				\item NIST Cybersecurity Framework (CSF)
				\item CISA Alerts
				\item Executive Order 14028
			\end{itemize}
		\end{block}
	\end{frame}
	
	\begin{frame}{AI as the Most Important Topic for Government Citations}
		\small
		\begin{block}{Why AI Research Gets the Most .gov Citations}
			\textbf{Why AI is valuable:}
			\begin{itemize}
				\item Executive Order 14110 (Safe, Secure, Trustworthy AI)
				\item NIST AI Risk Management Framework
				\item Federal AI research funding (NSF, DOE, DARPA)
				\item AI regulation through FTC, FCC, SEC, CFPB
			\end{itemize}
			\textbf{AI topics that get government attention:}
			\begin{itemize}
				\item AI safety and alignment
				\item Algorithmic bias and fairness
				\item Generative AI regulation (LLMs, deepfakes)
			\end{itemize}
		\end{block}
	\end{frame}
	
	
	% ============================================================
	% OSF (OPEN SCIENCE FRAMEWORK) INTEGRATION - SMALLER FONT
	% ============================================================
	
	\begin{frame}{OSF: Open Science Framework - Why It Matters for EB1A/EB2-NIW}
		\scriptsize
		\begin{block}{What is OSF?}
			\textbf{Answer:} OSF (osf.io) is a free, open-source research management platform maintained by the Center for Open Science (COS), which receives funding from federal agencies including the National Institutes of Health (NIH) and the National Science Foundation (NSF). OSF allows researchers to upload preprints, preprints, datasets, registered research protocols, and entire research workflows.
			\vspace{0.1cm}
			\textbf{Why OSF is strategically important for immigration evidence:}
			\begin{itemize}
				\item \textbf{OSF Preprints receive permanent DOIs} - your work becomes citable immediately, before journal peer review. Citations start accruing the day you upload.
				\item \textbf{OSF Preprints are harvested by SHARE} - a project of the U.S. Department of Education. Your work automatically becomes discoverable through .gov portals.
				\item \textbf{Federal agencies (NIH, NSF, DOE) mandate OSF for open data compliance} - using OSF demonstrates that your research meets federal standards for data management and transparency.
				\item \textbf{OSF Registrations create timestamped, immutable records} - proving you are the originator of specific research ideas. This is direct evidence of "original contributions" for EB1A.
				\item \textbf{OSF integrates with ORCID and Google Scholar} - citations to your OSF preprints appear in Google Scholar just like journal articles, and are tracked by major academic databases.
			\end{itemize}
		\end{block}
	\end{frame}
	
	\begin{frame}{OSF Preprints: The Connection to Government Recognition}
		\scriptsize
		\begin{block}{How OSF Preprints Lead to .gov and .edu Citations}
			\textbf{Answer:} When you upload a preprint to OSF, it is assigned a DOI (Digital Object Identifier) and becomes part of the global scholarly record. Unlike a personal blog or institutional repository, OSF Preprints are actively harvested by U.S. government systems.
			\vspace{0.1cm}
			\textbf{The government citation pathway:}
			\begin{itemize}
				\item \textbf{SHARE (share.osf.io)} is a joint project of the U.S. Department of Education and the Center for Open Science. SHARE crawls OSF Preprints and makes them searchable through .gov portals. Your preprint becomes part of the official U.S. government research database.
				\item \textbf{CHORUS (chorusaccess.org)} connects OSF Preprints to federal agency reports. When a federal agency publishes a report citing research, OSF preprints are included in the citation tracking.
				\item \textbf{PubMed Central (PMC)} accepts preprints that are later published as peer-reviewed articles. OSF preprints with DOIs can be referenced in NIH grant reports.
				\item \textbf{NASA ADS (Astrophysics Data System)} indexes OSF preprints in astronomy, physics, and related fields, making them searchable through Science.gov.
				\item \textbf{ERIC (Education Resources Information Center)} - the U.S. Department of Education's digital library - indexes OSF Preprints in education and workforce development through the EdArXiv community.
			\end{itemize}
			\textcolor{blue}{A single OSF preprint upload creates multiple pathways for your work to be discovered, downloaded, and cited by federal agency researchers and policymakers.}
		\end{block}
	\end{frame}
	
	\begin{frame}{OSF Preprint Communities: Aligning with Federal Priorities}
		\scriptsize
		\begin{block}{OSF Preprint Communities That Maximize Government Visibility}
			\textbf{Answer:} OSF organizes preprints into communities. Choosing the right community determines which federal agencies are most likely to discover and cite your work.
			\vspace{0.1cm}
			\textbf{High-impact communities for EB1A/EB2-NIW:}
			\begin{itemize}
				\item \textbf{CSRC Preprints} - the Computer Security Research Center community, aligned with NIST (National Institute of Standards and Technology). Research posted here is directly visible to NIST cybersecurity researchers who produce federal guidelines like the NIST Cybersecurity Framework. Citations from NIST reports become .gov citations.
				\item \textbf{EdArXiv} - the education preprint community, partnered with the U.S. Department of Education. Preprints posted here are automatically harvested by ERIC (.gov) and SHARE. This is the fastest pathway to .gov indexing for education and workforce development research.
				\item \textbf{SocArXiv} - social sciences preprints, harvested by SHARE and indexed in U.S. government social science databases. Relevant for policy research, economics, and workforce studies.
				\item \textbf{medRxiv and bioRxiv} - health and biomedical preprints. These are monitored by NIH researchers and can lead to citations in NIH grant reports, clinical practice guidelines, and federal health policy documents.
				\item \textbf{arXiv.org} - the original preprint server, now integrated with OSF. arXiv preprints are indexed by NASA ADS and appear in Science.gov. Particularly valuable for AI, machine learning, and cybersecurity research.
				\item \textbf{SSRN (Social Science Research Network)} - now part of Elsevier and integrated with OSF. SSRN preprints are indexed by the U.S. Commerce Research Library through EBSCO Business Source Complete, creating a pathway to Department of Commerce citations.
			\end{itemize}
			\textcolor{red}{The same research paper, uploaded to OSF through the right community, becomes discoverable through multiple .gov portals simultaneously.}
		\end{block}
	\end{frame}
	
	\begin{frame}{OSF Preprint DOIs: Why They Matter for Immigration}
		\scriptsize
		\begin{block}{DOIs (Digital Object Identifiers) and Government Recognition}
			\textbf{Answer:} Every OSF preprint receives a permanent DOI - a unique alphanumeric string that permanently identifies your work. DOIs are the standard for academic citation and are recognized by all major databases including Google Scholar, Crossref, Scopus, and Web of Science.
			\vspace{0.1cm}
			\textbf{Why OSF DOIs are valuable for EB1A/EB2-NIW:}
			\begin{itemize}
				\item \textbf{DOIs create a permanent, verifiable record of your contribution.} USCIS officers can click the DOI link and immediately see your preprint, its download counts, and its citation history. Unlike a personal website or blog, a DOI-backed preprint cannot be altered or disputed.
				\item \textbf{DOIs enable citation tracking across .gov domains.} When a federal agency report cites your OSF preprint, the citation includes the DOI. USCIS can verify that citation by resolving the DOI to your work. This creates a direct, unbroken chain of evidence from your research to government recognition.
				\item \textbf{DOIs are indexed by Google Scholar within 2-4 weeks.} Your OSF preprint will appear in Google Scholar search results alongside peer-reviewed journal articles. USCIS officers regularly check Google Scholar for citation evidence.
				\item \textbf{DOIs work with Crossref's Cited-by service.} When someone cites your OSF preprint, Crossref notifies you and updates the citation count. This automated tracking eliminates the need for manual citation discovery.
				\item \textbf{DOIs prove timestamped priority.} The DOI registration date is permanently recorded. If another researcher later publishes similar work, your DOI proves you were first. This directly supports the "original contributions" prong of EB1A.
			\end{itemize}
			\textcolor{blue}{A DOI is the academic equivalent of a patent filing date - it establishes priority and creates a permanent, verifiable record of your contribution.}
		\end{block}
	\end{frame}
	
	\begin{frame}{OSF Registrations: Timestamped Evidence of Original Contributions}
		\scriptsize
		\begin{block}{What OSF Registrations Are and Why They Matter}
			\textbf{Answer:} An OSF Registration is a timestamped, immutable snapshot of your research plan, methodology, or findings at a specific point in time. Unlike a preprint (which can be updated), a registration is frozen and cannot be altered after creation.
			\vspace{0.1cm}
			\textbf{Why registrations are powerful evidence for EB1A:}
			\begin{itemize}
				\item \textbf{Registrations prove you were the originator of specific research ideas.} The timestamp on a registration serves as legal evidence of priority. If you register a research protocol on OSF and later publish the results, the registration proves when you first conceived the study.
				\item \textbf{Federal agencies require pre-registration for funded research.} NIH and NSF both require or strongly encourage pre-registration of clinical trials and certain types of research. Complying with federal standards demonstrates that your research meets U.S. government quality expectations.
				\item \textbf{The Preregistration Challenge on OSF was funded by NSF.} The Center for Open Science received NSF funding to promote pre-registration. Participating in NSF-funded initiatives creates evidence of federal recognition.
				\item \textbf{Registrations appear in Google Scholar.} Like preprints, registrations are indexed and can be cited by other researchers, including federal agency researchers.
				\item \textbf{Registrations create a permanent, verifiable audit trail.} USCIS officers can click the registration DOI and see exactly what you proposed, when you proposed it, and whether you executed it as planned. This transparency supports claims of scholarly integrity and contribution.
			\end{itemize}
			\textcolor{red}{An OSF Registration with a DOI is admissible evidence of "original contribution of major significance" because it establishes timestamped priority of research ideas.}
		\end{block}
	\end{frame}
	
	\begin{frame}{OSF and Federal Data Management Plans (DMPs)}
		\scriptsize
		\begin{block}{Meeting Federal Data Management Requirements Through OSF}
			\textbf{Answer:} Most U.S. federal research agencies now require Data Management Plans (DMPs) for grant applications. OSF integrates directly with DMPTool, the federally approved DMP creation platform.
			\vspace{0.1cm}
			\textbf{Which agencies require DMPs that OSF supports:}
			\begin{itemize}
				\item \textbf{National Science Foundation (NSF)} - requires a two-page DMP for all grant proposals. DMPTool is the NSF-recommended platform. OSF projects can be linked to DMPs created in DMPTool.
				\item \textbf{National Institutes of Health (NIH)} - requires DMPs for all grant proposals requesting \$500,000 or more in direct costs. OSF meets NIH data sharing requirements.
				\item \textbf{Department of Energy (DOE)} - requires DMPs for all research funded by DOE. OSF is approved for DOE data management.
				\item \textbf{Department of Defense (DoD)} - requires DMPs for basic research funding. OSF meets DoD requirements for unclassified research.
			\end{itemize}
			\textbf{Why DMP compliance matters for immigration:}
			\begin{itemize}
				\item Demonstrates that your research meets federal quality and transparency standards
				\item Shows you understand and follow U.S. government research requirements
				\item Creates additional documentation for your EB2-NIW "national importance" argument
				\item Links your work to specific federal agency funding priorities
			\end{itemize}
		\end{block}
	\end{frame}
	
	\begin{frame}{OSF Metrics: Quantifying Your Research Impact}
		\scriptsize
		\begin{block}{What Metrics OSF Provides and How to Use Them}
			\textbf{Answer:} OSF automatically tracks and displays several metrics for every preprint and project. These metrics provide quantifiable evidence of your research impact.
			\vspace{0.1cm}
			\textbf{Metrics OSF tracks:}
			\begin{itemize}
				\item \textbf{Downloads} - the number of times your preprint has been downloaded. High download counts indicate that other researchers and practitioners find your work valuable. International downloads demonstrate global reach.
				\item \textbf{Views} - the number of times your OSF project page has been viewed. This includes visits from .edu and .gov domains (which you can document).
				\item \textbf{Citations} - OSF integrates with Google Scholar to display citation counts directly on your preprint page. USCIS officers can see at a glance how many times your work has been cited.
				\item \textbf{Fork count} - how many other researchers have created copies (forks) of your OSF project to reuse your methodology. Forks prove that other researchers are building directly on your work.
				\item \textbf{Storage usage} - less relevant, but shows the scale and complexity of your research.
			\end{itemize}
			\textbf{How to use OSF metrics in your immigration petition:}
			\begin{itemize}
				\item Screenshot your OSF preprint page showing download counts in the thousands
				\item Include metrics in your "sustained acclaim" evidence for EB1A
				\item Compare your metrics to field averages to demonstrate above-normal impact
				\item Track increases over time to show "sustained" rather than one-time impact
			\end{itemize}
		\end{block}
	\end{frame}
	
	\begin{frame}{OSF + Regulations.gov: The Complete Policy Influence Pathway}
		\scriptsize
		\begin{block}{How OSF Preprints Strengthen Policy Comments}
			\textbf{Answer:} The most powerful combination for EB1A "major significance" evidence is an OSF preprint that you then cite in a Regulations.gov policy comment.
			\vspace{0.1cm}
			\textbf{The complete pathway:}
			\begin{enumerate}
				\item You publish a research preprint on OSF, receiving a permanent DOI.
				\item Your OSF preprint is indexed by Google Scholar, SHARE, and other government systems.
				\item You submit a comment to a federal agency's Request for Information (RFI) or Notice of Proposed Rulemaking (NPRM) on Regulations.gov.
				\item In your comment, you cite your own OSF preprint DOI as evidence supporting your recommendations.
				\item The agency considers your comment as part of the official record. If the agency cites your comment or your preprint in the final rule, you receive a .gov citation.
			\end{enumerate}
			\textbf{Why this pathway is uniquely powerful:}
			\begin{itemize}
				\item Your OSF preprint provides the scholarly evidence. Your Regulations.gov comment provides the policy application. Together, they prove your work has both academic merit AND real-world policy influence.
				\item The DOI creates a verifiable link between your research and the policy comment. USCIS can click the DOI and immediately see your preprint.
				\item This addresses both EB1A (original contributions) and EB2-NIW (national importance) simultaneously.
			\end{itemize}
			\textcolor{red}{A single research project, properly executed through OSF and Regulations.gov, can generate evidence for multiple immigration criteria.}
		\end{block}
	\end{frame}
	
	\begin{frame}{OSF: Summary of Immigration Benefits}
		\scriptsize
		\begin{block}{Why Every EB1A/EB2-NIW Applicant Should Use OSF}
			\textbf{Answer:} OSF provides unique advantages that cannot be obtained through traditional journal publication alone.
			\vspace{0.1cm}
			\textbf{Eight reasons OSF strengthens your immigration case:}
			\begin{enumerate}
				\item \textbf{Permanent DOIs} - your work becomes citable immediately, not after months of peer review.
				\item \textbf{.gov harvesting} - SHARE (Department of Education) automatically indexes your work.
				\item \textbf{Federal compliance} - OSF meets NIH, NSF, and DOE data management requirements.
				\item \textbf{Timestamped registrations} - proves you originated specific research ideas.
				\item \textbf{Download metrics} - quantifiable evidence of impact and sustained acclaim.
				\item \textbf{Integration with Regulations.gov} - cite your own preprints in policy comments.
				\item \textbf{Google Scholar indexing} - your OSF preprints appear alongside journal articles.
				\item \textbf{Free and accessible} - no cost barrier to building your evidence portfolio.
			\end{enumerate}
			\textcolor{blue}{Traditional journal publication can take 6-18 months. OSF preprints are citable within hours. For immigration purposes, starting your citation clock earlier is always better.}
		\end{block}
	\end{frame}
	
	% ============================================================
	% SECTION 10: CONGRESSIONAL AND SENATOR SUPPORT LETTERS
	% ============================================================
	
	\begin{frame}{How to Get a Congressman or Senator Support Letter}
		\small
		\begin{block}{Step-by-Step Strategy}
			\textbf{Step 1 - Identify the right office:}
			\begin{itemize}
				\item Your local Representative
				\item Senators from your state
				\item Members with expertise in your field
			\end{itemize}
			\textbf{Step 2 - Prepare request packet:}
			\begin{itemize}
				\item Cover letter explaining national importance
				\item 2-3 page research summary
				\item Citation report (.gov, .edu, journal)
				\item Draft letter of support
			\end{itemize}
			\textbf{Step 3 - Submit:} Website contact form or mail to district office.
		\end{block}
	\end{frame}
	
	\begin{frame}{What to Include in a Congressional Support Letter Request}
		\small
		\begin{block}{Components of a Successful Request}
			\textbf{Draft letter should include:}
			\begin{itemize}
				\item Statement of your expertise
				\item Description of research and applications
				\item How your work serves U.S. national interests
				\item Evidence of recognition (.gov citations, peer review)
				\item Specific USCIS wording ("extraordinary ability," "national importance")
			\end{itemize}
			\textbf{What NOT to ask:} Do not ask them to endorse immigration decisions or contact USCIS directly.
		\end{block}
	\end{frame}
	
	% ============================================================
	% SECTION 11: CONCLUSION AND FINAL CHECKLIST
	% ============================================================
	
	\begin{frame}{Complete EB1A/EB2-NIW Citation Evidence Checklist}
		\small
		\begin{block}{Before Submitting Your Petition}
			\textbf{Domain 1: Government Citations (.gov)}
			\begin{itemize}
				\item [$\square$] Regulations.gov comments (with docket numbers)
				\item [$\square$] Federal agency citations (BLS, Fed, DOE, DOL)
				\item [$\square$] CLEAR DOL dataset indexing
				\item [$\square$] Science.gov indexing confirmation
			\end{itemize}
			\textbf{Domain 2: Educational Impact (.edu)}
			\begin{itemize}
				\item [$\square$] ERIC-indexed publications
				\item [$\square$] University citations (Google Scholar .edu)
				\item [$\square$] Course material inclusions
			\end{itemize}
		\end{block}
	\end{frame}
	
	\begin{frame}{Complete EB1A/EB2-NIW Checklist (continued)}
		\small
		\begin{block}{Before Submitting Your Petition}
			\textbf{Domain 3: Peer Review \& Expert Recognition}
			\begin{itemize}
				\item [$\square$] 5+ peer review certificates (minimum)
				\item [$\square$] Editorial board memberships
				\item [$\square$] Invitation emails from journals
			\end{itemize}
			\textbf{Domain 4: Publication \& Citation Metrics}
			\begin{itemize}
				\item [$\square$] Google Scholar profile with citations
				\item [$\square$] Scopus/Web of Science citation report
				\item [$\square$] Total citations (aim for 50+)
			\end{itemize}
			\textbf{Domain 5: Supporting Evidence}
			\begin{itemize}
				\item [$\square$] Congressional or Senator support letter
				\item [$\square$] Expert opinion letters
			\end{itemize}
		\end{block}
	\end{frame}
	
	\begin{frame}{How to Present Your Evidence to USCIS}
		\small
		\begin{block}{Documentation Strategy for Maximum Impact}
			\textbf{Structure your exhibit tabs:}
			\begin{enumerate}
				\item \textbf{Tab A}: .gov Citations
				\item \textbf{Tab B}: .edu Citations
				\item \textbf{Tab C}: Journal Citations
				\item \textbf{Tab D}: Peer Review Evidence
				\item \textbf{Tab E}: Support Letters
				\item \textbf{Tab F}: Portfolio of Publications
			\end{enumerate}
			\textbf{For each citation:} Screenshot, highlight your name, include URL, save as PDF.
		\end{block}
	\end{frame}
	
	\begin{frame}{Timeline for Building Your Citation Portfolio}
		\small
		\begin{block}{Realistic Timeline from Zero to EB1A-Ready (12-24 months)}
			\textbf{Months 1-3:} Set up Zotero, Google Scholar; identify journals; draft first paper.
			\vspace{0.1cm}
			\textbf{Months 4-8:} Submit to journal; begin Regulations.gov commenting; apply for peer review.
			\vspace{0.1cm}
			\textbf{Months 9-12:} Track citations; submit second paper; request congressional letter.
			\vspace{0.1cm}
			\textbf{Months 13-18:} Upload to multiple platforms; seek .edu and government citations.
			\vspace{0.1cm}
			\textbf{Months 19-24:} Organize exhibits; draft petition; consult attorney.
		\end{block}
	\end{frame}
	
	\begin{frame}{Common Mistakes to Avoid}
		\small
		\begin{block}{What Can Hurt Your Immigration Case}
			\textbf{Evidence mistakes:}
			\begin{itemize}
				\item Submitting predatory journal publications
				\item Claiming unverifiable citations
				\item Missing .gov or .edu evidence
			\end{itemize}
			\textbf{Strategy mistakes:}
			\begin{itemize}
				\item Focusing only on volume (quality matters more)
				\item Ignoring policy and government engagement
				\item Neglecting documentation
			\end{itemize}
			\textbf{Petition mistakes:}
			\begin{itemize}
				\item Poorly organized exhibits
				\item Missing URLs for online evidence
			\end{itemize}
			\textcolor{red}{An RFE is common - prepare by keeping all documentation.}
		\end{block}
	\end{frame}
	
	\begin{frame}{Final Success Checklist}
		\small
		\begin{block}{Before You Submit Your Immigration Petition}
			\begin{itemize}
				\item [$\square$] All .gov citations have working URLs
				\item [$\square$] .edu citations from accredited U.S. universities
				\item [$\square$] Peer review certificates signed and dated
				\item [$\square$] Google Scholar profile public and updated
				\item [$\square$] Publication PDFs clean and readable
				\item [$\square$] Exhibits organized and numbered
				\item [$\square$] Consulted with qualified immigration attorney
				\item [$\square$] Prepared for potential RFE with backup evidence
			\end{itemize}
			\textcolor{blue}{USCIS approves thousands of EB1A and EB2-NIW cases each year.}
		\end{block}
	\end{frame}
	
	\begin{frame}{Resources and Next Steps}
		\small
		\begin{block}{Where to Go From Here}
			\textbf{Immediate actions (next 7 days):}
			\begin{enumerate}
				\item Set up Google Scholar profile (free)
				\item Install Zotero (free)
				\item Identify 2-3 Regulations.gov RFIs
				\item Create list of 5 target journals
			\end{enumerate}
			\textbf{Key websites to bookmark:}
			\begin{itemize}
				\item regulations.gov
				\item science.gov
				\item eric.ed.gov
				\item clear.dol.gov
				\item pubmed.ncbi.nlm.nih.gov
			\end{itemize}
		\end{block}
	\end{frame}
	
	\begin{frame}{Disclaimer and Legal Advice}
		\small
		\begin{block}{Important Notice}
			\textbf{Disclaimer:} This provides educational guidance only and does not constitute legal advice. Immigration laws change regularly.
			\vspace{0.2cm}
			\textbf{Recommendation:}
			\begin{itemize}
				\item Always consult a qualified immigration attorney
				\item Each case is unique
				\item USCIS officers have discretion
				\item This teaches strategies, not guarantees
			\end{itemize}
			\textcolor{blue}{Use this course to prepare evidence, then work with an attorney.}
		\end{block}
	\end{frame}
	
	\begin{frame}{Thank You and Good Luck!}
		\centering
		\vspace{0.5cm}
		\textcolor{blue}{\Large Your Path to EB1A/EB2-NIW Success Starts Here}
		\vspace{0.3cm}
		\small
		\begin{itemize}
			\item Implement strategies consistently
			\item Track your progress monthly
			\item Document everything systematically
			\item Build relationships with experts
			\item Stay persistent - impact takes time
		\end{itemize}
		\vspace{0.3cm}
		\textcolor{red}{\textbf{The best time to start was yesterday.}}
		\textcolor{red}{\textbf{The second best time is now.}}
		\vspace{0.3cm}
		\textcolor{blue}{\Large Good luck on your immigration journey!}
	\end{frame}
	
\end{document}